% steps/step11.tex
\section[گام یازدهم: استقرار سیستم]{\faCogs\ \ گام یازدهم: استقرار سیستم}
\begin{stepbox}

\field{انتخاب بستر مناسب استقرار (On-Premise Deployment)}{
برای کاربردی شدن مدل در محیط بالینی، باید آن را از فضای آزمایشگاهی (مانند نوت‌بوک‌های پایتون) خارج کرده و بدون ایجاد مزاحمت، در جریان کار روزمره پزشک قرار داد. در حوزه پزشکی به دلیل حساسیت بالای داده‌ها و قوانین سخت‌گیرانه حریم خصوصی، ارسال تصاویر بیماران به سرویس‌های ابری (\lr{Cloud}) معمولاً منطقی یا حتی مجاز نیست. بنابراین، مدلِ بهینه‌شده روی سرورهای داخلی و امن بیمارستان (\lr{On-Premise}) یا روی زیرساخت نزدیک به دستگاه تصویربرداری (\lr{Edge/On-site Server}) مستقر می‌شود تا کنترل کامل امنیت، دسترسی و نگه‌داری داده‌ها حفظ گردد.
}

\field{یکپارچه‌سازی با سیستم‌های موجود بیمارستان (Integration)}{
رادیولوژیست‌ها قرار نیست با کد، محیط‌های ناشناخته یا ابزارهای توسعه کار کنند؛ آن‌ها از سامانه‌های استاندارد مانند \lr{PACS} (سیستم بایگانی و تبادل تصاویر پزشکی) برای مشاهده و گزارش‌نویسی استفاده می‌کنند. در نتیجه، سیستم سگمنتیشن باید به‌صورت یک جزء پنهان و قابل اطمینان به \lr{PACS} متصل شود و خروجی را در قالب استانداردهای پزشکی (مانند \lr{DICOM} و \lr{DICOM-SEG}/\lr{RTSTRUCT} در صورت نیاز) تحویل دهد. این یکپارچه‌سازی تضمین می‌کند که پزشک خروجی مدل را در همان محیط کاری معمول خود و به‌صورت یک لایه \lr{Overlay} روی تصویر مشاهده کند.
}

\field{توسعه API و بسته‌بندی سرویس (API \& Dockerization)}{
برای ارائه مدل به‌عنوان یک سرویس پایدار، تکرارپذیر و قابل استقرار در شبکه بیمارستان، مدل نهایی به همراه تمام پیش‌نیازها (پیش‌پردازش، پس‌پردازش و وابستگی‌ها) در قالب یک کانتینر مانند \lr{Docker} بسته‌بندی می‌شود. سپس از طریق یک وب‌سرویس (برای مثال با \lr{FastAPI}) یک \lr{API} داخلی فراهم می‌گردد تا سامانه‌های بیمارستانی بتوانند تصاویر را ارسال کرده و خروجی سگمنتیشن را دریافت کنند. این رویکرد، استقرار و به‌روزرسانی نسخه‌های مختلف مدل را ساده‌تر و کنترل‌پذیرتر می‌کند.
}

\field{جریان کار عملیاتی (Workflow) در دنیای واقعی}{
در سناریوی واقعی، روند پیشنهادی به شکل زیر است:  
\lr{CT Scan} بیمار گرفته می‌شود $\rightarrow$ تصویر به‌صورت خودکار به \lr{API} مدل ارسال می‌گردد $\rightarrow$ مدل در زمان کوتاه (نزدیک به بلادرنگ) کبد/تومور را سگمنت می‌کند $\rightarrow$ ماسک خروجی (یا ساختار سگمنتیشن استاندارد) به \lr{PACS} بازگردانده می‌شود $\rightarrow$ پزشک هنگام مشاهده پرونده بیمار، ناحیه پیشنهادی مدل را به‌صورت یک \lr{Overlay} رنگی روی تصویر می‌بیند و تصمیم نهایی و مسئولیت بالینی را بر عهده می‌گیرد.
}

\end{stepbox}
